(1) 相对论动量 $p=\gamma mv$，洛伦兹力在一维为 $dp/dt=qE\equiv F$。初态静止给 $p(t)=Ft$；能量—动量关系 $W^2=m^2c^4+p^2c^2$ 得 $v(t)=pc^2/W=Ft/[m\sqrt{1+(Ft/(mc))^2}]$，严格满足 $|v|<c$。
(2) 定义 $a_0=F/m$，积分 $v(t)=a_0t/\sqrt{1+(a_0t/c)^2}$，得 $x(t)-x(0)=\frac{c^2}{a_0}[\sqrt{1+(a_0t/c)^2}-1]$；在 $t=0$ 位移为零，求导恰回到速度公式，不存在超光速位移斜率。
(3) 相对论总能量 $W(t)=mc^2\sqrt{1+(a_0t/c)^2}$，动能 $K(t)=W-mc^2=mc^2[\sqrt{1+(a_0t/c)^2}-1]$。将位移乘以 $qE=ma_0$，得 $qE[x-x_0]=mc^2[\sqrt{1+(a_0t/c)^2}-1]=K(t)$，能量由电源做功提供。
(4) 当 $a_0t/c\ll1$，展开得 $v\simeq a_0t$、$x-x_0\simeq a_0t^2/2$、$K\simeq ma_0^2t^2/2$。当 $t\to\infty$，$v\to c^-$，$x-x_0=ct-c^2/a_0+o(1)$；真正带电粒子在外场中加速时通常辐射，强场或长时间需评价辐射能耗及自力，题设无辐射是假设。
